\documentclass{situationist-preprint}

\title{Berkeley's Idealism and the Cost of Primitives}
\msid{MET-03-b}
\msfamily{MET-03}
\msdiscipline{Metaphysics}
\mstype{Consequence}
\mskeywords{Berkeley; idealism; ideology; parsimony; material substance; notions; master argument}

\begin{document}
\maketitle

\begin{abstract}
MET-03-a argues that idealism is consistent but needs one more primitive than realism, the notion of conception, and is eliminated by a principle of parsimony of ideology. We draw three consequences for Berkeley's own case. First, his argument that matter is idle succeeds against Lockean material substance but not against the structure that every adequate theory needs, which Berkeley himself retains as the archetypes in God's mind. Second, his doctrine that we have notions, not ideas, of spirits concedes that conception is not a species of experience, which is the premise on which the extra primitive rests. Third, the master argument fails, as the consistency of idealism's rival requires. Berkeley's case against matter, costed by his own standard, counts against him.
\end{abstract}

\section{Berkeley's argument from idleness}

MET-03-a compares realism and idealism by their ideology, the notions each takes as primitive. Every adequate theory of the order of experience needs experience, structure and production (MET-03-a, Assumption~2.3); realism needs nothing more (Definition~2.4); idealism needs, in addition, conception, the relation in which a mind stands to a state it does not perceive by sense (Definition~2.5 and Assumption~2.6). Lemma~3.2 counts the difference, and Theorem~3.3 eliminates idealism by the parsimony principle P1.

Berkeley argued for idealism by an appeal to economy of his own. Even if there were solid, figured, movable substances without the mind, he wrote, we could not know it, since the same ideas could be produced in us without them, as is granted in dreams; and if there were no such substances, we might have the very same reasons to think there were that we have now \citep[\S\S 18--20]{BerkeleyPrinciples}. Matter, on this argument, is idle: God could produce the ideas directly, and the hypothesis of matter adds an entity that explains nothing.

The argument has two targets that Berkeley did not distinguish. The first is Lockean material substance: a substratum that supports qualities, that is neither an idea nor a structure of relations among ideas, and whose only role is to be the bearer of qualities. The second is structure in the sense of MET-03-a (Definition~2.2): a system of states with lawlike relations, not identical to any pattern among experiences, which the order of experience requires for its counterfactuals about places where no observer is.

\begin{proposition}\label{p:substratum}
Berkeley's argument from idleness eliminates material substratum and does not eliminate structure.
\end{proposition}

\begin{proof}
A substratum, in Locke's sense, adds a primitive to any adequate theory, the relation of supporting qualities, without adding explanatory power, since the order of experience is explained by structure and production alone (MET-03-a, Assumption~2.7). A theory with a substratum is therefore eliminated by P1 of MET-03-a in favor of the same theory without it, which is the conclusion Berkeley wanted. Structure, by contrast, belongs to the ideology of every adequate theory (Assumption~2.3), Berkeley's included: the archetypes in the mind of God, by which he secures the continued existence of the unperceived tree \citep[third dialogue]{BerkeleyDialogues}, and the \enquote{set rules or established methods} by which God excites ideas in us \citep[\S 30]{BerkeleyPrinciples}, are a structure in the sense of Definition~2.2. An argument from idleness cannot eliminate a notion that the arguer's own theory needs.
\end{proof}

Proposition~\ref{p:substratum} shows what Berkeley's argument really establishes. It is an argument for a theory whose ideology consists of experience, structure and production, without substratum. That theory is realism as MET-03-a defines it. Berkeley's further step, from the idleness of substratum to the dependence of structure on a conceiving mind, is not supported by idleness at all; it adds a primitive rather than removing one.

\section{Notions and the indefinability of conception}

MET-03-a rests the extra cost of idealism on Assumption~2.6: a theory that makes structure depend on being conceived has conception among its primitives, because conception cannot be defined as experience, as production, or in terms of structure. A Berkeleyan might hope to define it as experience, since for Berkeley to be is to be perceived \citep[\S 3]{BerkeleyPrinciples}, and perception is experience.

Berkeley closed that route himself. God perceives nothing by sense, and has no sensations \citep[third dialogue]{BerkeleyDialogues}; and in the edition of 1734 Berkeley added that we have no idea of spirits, of their acts, or of relations, but only a \emph{notion} of them \citep[\S\S 27, 89, 142]{BerkeleyPrinciples}. A notion, in his sense, is a way of understanding what is not and cannot be an idea. God's relation to the archetypes is an act of a spirit, and so is known by notion, not by idea; it is not a sensory experience in any mind. This is Assumption~2.6 in Berkeley's own vocabulary. The doctrine of notions, which Berkeley needed to answer the objection that his principles left him without any knowledge of minds, is precisely the admission that his theory employs a notion that is not definable in terms of experience \citep{Winkler1989}.

\begin{corollary}\label{c:berkeley}
By Berkeley's own principles, the ideology of his theory includes a notion not definable in terms of experience, structure or production, and his theory is eliminated by parsimony.
\end{corollary}

\begin{proof}
The relation of God to the archetypes is not an idea and not an experience, by the doctrine of notions and the denial that God has sensations; it is not production, since God's conceiving an archetype is not his producing an idea in a finite mind; and it is not structure, since it is the relation on which, for Berkeley, structure depends. So it is a further primitive, and Theorem~3.4 of MET-03-a applies.
\end{proof}

\section{The master argument}

Berkeley also offered an argument meant to show that the realist's thesis is not merely idle but incoherent: try to conceive a tree existing unconceived, and you will find that you are conceiving it \citep[\S 23]{BerkeleyPrinciples}. If the argument succeeded, realism would not be adequate, Assumption~2.7 of MET-03-a would fail, and the parsimony comparison would never begin.

The argument fails, and MET-03-a gives the reason in a form Berkeley could not have resisted. Realism's distinctive thesis, that some structural states produce no experience in any mind, is stated in the shared vocabulary and has a model; the model of Proposition~3.1 of MET-03-a, with the conceiving mind removed, is a model of it. What the master argument shows is that whoever conceives that a tree is unconceived conceives the tree; it does not show that the tree is conceived, since the content conceived is that the tree is not. The distinction between conceiving something and conceiving that it exists unconceived is the one on which \citet{Prior1955} and \citet{Pitcher1977} rest their diagnoses, and it is enough to keep realism among the adequate theories.

\section{Conclusion}

Berkeley's idealism is consistent, and his case against matter rests on a sound principle. Applied with care, the principle removes Lockean substratum and leaves structure, which Berkeley needs and keeps under another name. His doctrine of notions concedes that the relation by which God sustains that structure is a further primitive, and his master argument does not show that realism is incoherent. By the standard of economy that Berkeley himself invoked, the theory that remains is realism.

\begin{thebibliography}{9}\small
\bibitem[Berkeley(1710)]{BerkeleyPrinciples} Berkeley, G. (1710). \emph{A Treatise Concerning the Principles of Human Knowledge}. Dublin. 2nd edition, London, 1734.
\bibitem[Berkeley(1713)]{BerkeleyDialogues} Berkeley, G. (1713). \emph{Three Dialogues between Hylas and Philonous}. London.
\bibitem[Pitcher(1977)]{Pitcher1977} Pitcher, G. (1977). \emph{Berkeley}. Routledge \& Kegan Paul.
\bibitem[Prior(1955)]{Prior1955} Prior, A.~N. (1955). Berkeley in logical form. \emph{Theoria}, 21, 117--122.
\bibitem[Winkler(1989)]{Winkler1989} Winkler, K.~P. (1989). \emph{Berkeley: An Interpretation}. Clarendon Press.
\end{thebibliography}

\end{document}
